% ============================================
% Johns Hopkins Letter Template
% Uses the built-in 'letter' class
% ============================================

\documentclass[11pt,letterpaper]{letter}

\usepackage{graphicx}
\usepackage{eso-pic}
\usepackage{geometry}
\usepackage{microtype}
\usepackage[utf8]{inputenc}
\usepackage[hidelinks]{hyperref}

% ----- Page Setup -----
\geometry{left=25mm,right=25mm,top=25mm,bottom=25mm}

% ----- Logo Placement (foreground, first page only) -----
\newcommand{\PutJHULogoFG}[1]{%
  \AddToShipoutPictureFG*{%
    \AtPageUpperLeft{%
      \hspace{10mm}%
      \raisebox{-38mm}{%
        \includegraphics[width=6cm]{#1}%
      }%
    }%
  }%
}

% ----- Sender Information -----
\signature{Decory Edwards}
\address{Department of Economics \\ Johns Hopkins University \\ Baltimore, MD 21218 \\ \texttt{dedwar65@jh.edu}}

% ----- Recipient Info -----
\begin{document}
\PutJHULogoFG{Images/jhulogo.png}

\begin{letter}{Hiring Committee \\
Board of Governors of the Federal Reserve System \\
Washington, DC}

\opening{Dear Hiring Committee,}

I am writing to express my interest in the Economist position at the Board of Governors of the Federal Reserve System. I am a Ph.D.\ candidate in the Department of Economics at Johns Hopkins University (advisors: Christopher D.\ Carroll, Nicholas W.\ Papageorge, and Stelios Fourakis), and I expect to receive my degree in May 2026. My primary specializations are macroeconomic modeling, wealth and income dynamics, and computational economics.

My research agenda focuses on the ability of macroeconomic models to generate empirically realistic distributions of wealth and income over the life cycle. A key driver of that work is modeling heterogeneity in the rate of return on assets, and understanding how return dispersion contributes to portfolio accumulation across households.

In my job market paper, I calibrate a life cycle heterogeneous agent model to match wealth moments from the Survey of Consumer Finances by allowing households to differ in their portfolio returns. The model helps explain observed wealth concentration and return dispersion. In a related research project using the Health and Retirement Study, I examine how trust influences both income growth and asset returns. I find a hump shaped relationship for returns and characterize the distribution of the persistent component of returns to net worth. These experiences have equipped me to frame research strategies and design modeling approaches, execute econometric and simulation analysis with large micro data sets, and translate results into clear and rigorous written deliverables for both academic and practitioner audiences.

I am particularly interested in this position because it combines current economic analysis, forecasting, and longer run research in direct support of monetary policy and financial regulation. Much of the Board’s work requires careful interpretation of incoming data, development and maintenance of econometric models, and clear communication of uncertainty and alternative scenarios to senior staff and policymakers. For example, forecasting exercises and Tealbook contributions require judgment about which movements in aggregate data reflect transitory shocks versus persistent structural change. My training has emphasized precisely these issues. In my research, I routinely distinguish persistent from transitory components, assess model fit against multiple empirical moments, and evaluate how alternative assumptions affect policy relevant conclusions. This background aligns closely with the analytical demands of policy briefing, forecasting, and model development at the Board.

I am enthusiastic about the opportunity to live and work in Washington, DC and to engage fully in the in person collaborative environment that characterizes the Federal Reserve Board. I am comfortable working on site according to the Board’s expectations and value the close interaction with economists, analysts, and policymakers that this enables. I am also drawn to the prospect of a sustained multi year role, which would allow me to contribute meaningfully to policy analysis while continuing to develop my research agenda in an environment that values both rigorous applied work and original scholarship.

I believe I would be a strong fit for the Federal Reserve Board because:
\begin{enumerate}
    \item My experience with structural modeling and empirical validation aligns with the Board’s use of econometric models for forecasting and policy analysis.
    \item My research on heterogeneity prepares me to analyze distributional implications of macroeconomic and regulatory policies.
    \item I bring strong written and verbal communication skills that are essential for producing policy memoranda, briefing materials, and research outputs for diverse audiences.
\end{enumerate}

I would welcome the opportunity to contribute my quantitative training, research experience, and interest in policy relevant analysis to the work of the Federal Reserve Board. Thank you for considering my application.

\closing{Sincerely,}

\end{letter}
\end{document}
